\documentclass[12pt]{article}
\usepackage[utf8]{inputenc}
\usepackage[T1]{fontenc}
\usepackage{hyperref}
\title{Kulisusu Language Resources for K-12 Curriculum Development}
\author{Compiled by Claude Sonnet 5 (claude-sonnet-5) for Andr\'as Kornai}
\date{}
\begin{document}
\maketitle

\section{Introduction}
Kulisusu is an Austronesian (Celebic, Bungku--Tolaki) language of northern Buton Island, Southeast Sulawesi, centred on North Buton Regency. Wikipedia cites about 23,000 speakers. It forms a dialect chain with Koroni and Taloki. It is written in Latin script, but no standard spelling system is reported.

\section{Primary Educational Resources}
\begin{itemize}
\item ``Formation of Social Values through the Use of Kulisusu Language in North Buton Regency Elementary School Students'' (study of Kulisusu use in primary schools): \url{https://ejurnal.warmadewa.ac.id/index.php/jret/article/view/8100}
\end{itemize}
No textbooks or primers were found.

\section{Linguistic and Reference Materials}
\begin{itemize}
\item Kulisusu language (Wikipedia): \url{https://en.wikipedia.org/wiki/Kulisusu_language}
\item Omniglot page on Kulisusu: \url{https://www.omniglot.com/writing/kulisusu.htm}
\item Verbix language page: \url{https://docs.verbix.com/Languages/Kulisusu}
\item Kulisusu language maintenance among teenagers in Ereke (sociolinguistic study): \url{https://ejurnal.warmadewa.ac.id/index.php/jret/article/view/323}
\item Article on Etnosia (Universitas Hasanuddin): \url{https://journal.unhas.ac.id/index.php/etnosia/article/view/2576/pdf}
\item Article in Jurnal Ilmiah Universitas Dayan: \url{https://ejournal.unidayan.ac.id/index.php/JIU/article/download/2260/369}
\end{itemize}
I did not open the last two articles, so their contents are unverified. The search summary links one of the Indonesian studies to an Android dictionary app (3,100 words, Indonesian to Kulisusu). I could not tell which of these links that study is.

\section{Religious and Community Resources}
\begin{itemize}
\item Praying for Indonesia, ``Who are the Kulisusu?'': \url{https://www.prayingforindonesia.com/ethnic-groups/the-people-of-sulawesi/who-are-the-kulisusu/}
\end{itemize}

\section{Digital and Audio Resources}
No audio or video resources were found. The Android dictionary app mentioned above was not located as a download.

\section{Curriculum Development Context}
Kulisusu is stable and spoken by the whole community, but not known to be taught in schools. The local studies suggest some interest in using it in primary schools and in maintaining it among young people. Next steps are to standardise an orthography, build on the existing wordlist or dictionary work, and explore local-content school provisions with North Buton education authorities.
\end{document}